\begin{proposition}[an explicit eight-dimensional coboundary]\label{push:src:C.3.25}
Let $\varphi$ be an orientation-preserving equivariant return map,
normalized to be the identity on a fixed tubular neighbourhood.
Let $V_\varphi$ be its actual framed-surgery six-manifold, constructed
as in \ref{push:src:C.3.14}. Choose the sign of $x$ so that the unit circle
bundle of its complex line, with the boundary fibre orientation,
is the actual oriented free circle bundle. Orient its meridian
accordingly, so that $\langle x,S^2\rangle=1$. This fixes the
possible associated-line versus dual-line convention. Write
\[
 p=\langle p_1(V_\varphi)x,[V_\varphi]\rangle,
 \qquad k=\langle x^3,[V_\varphi]\rangle.
\]
There is a specified smooth spin$^c$ coboundary $Z_\varphi$ of
the actual clutching sphere $\Sigma_\varphi$. Its fourth
intersection form and characteristic covectors are
\[
 Q=\begin{pmatrix}k&1\\1&0\end{pmatrix},\qquad
 p_1(Z_\varphi)=(p+k,3),\qquad c^2=(k,1),
\]
where $c$ is the determinant class and the covectors are evaluated
on the ordered integral basis specified below. In particular
\[
 \sigma(Z_\varphi)=0,\qquad
 p_1^2=6p-3k,\qquad p_1c^2=p+k,\qquad c^4=k.
\]
With the boundary-consistent orientations, its spin$^c$ index
formula gives
\[
 \mu(\Sigma_\varphi)\equiv\frac{p-4k}{672}\pmod\Z.
\]
Boundary connected sum with $S^1\times D^7$ gives a specified
smooth spin$^c$ coboundary of the actual mapping torus, with the
same displayed characteristic numbers. This construction makes
no use of Wall's knot coordinate or the Davis generator comparison.
\end{proposition}
\Uses{the hypotheses, notation and definitions stated here; \cref{push:src:C.3.14}; \cref{push:src:C.2.1}; \citet[Section 2.2, equation (2.1)]{CrowleyNordstrom2020}; \citet[Section 6, equation (11)]{EellsKuiper1962}; \cref{push:src:C.1.1}}

\begin{proof}
We first identify the boundary geometrically. Write $C$ for the
exterior of the actual fixed three-knot in the orbit six-sphere.
The framed-surgery construction is
\[
 V_\varphi=C\cup_{S^3\times S^2}(D^4\times S^2).
\]
This description, including its framing, is the direct cut-and-paste
construction of \ref{push:src:C.3.14}; its later use of knot classification is not
needed here. It is spin and has $H^2=H^4=\Z$, $H^3=0$.
Let $L_x\to V_\varphi$ be the complex line with first Chern class
$x$, and put $Z_0=D(L_x)$, oriented by the base followed by its
complex line fibre. Its boundary $T=S(L_x)$ is the corresponding
principal circle bundle.

On the surgery piece $N=S^2\times D^4$, the line is the product
of the Euler-one line $\mathcal O(1)\to S^2$ and $D^4$.
Thus $T$ contains the specified product $S^3_{\rm Hopf}\times D^4$.
Attach one four-handle
\[
 D^4_{\rm Hopf}\times D^4
 \quad\hbox{along}\quad S^3_{\rm Hopf}\times D^4
\]
using this product identification and framing. Denote the resulting
eight-manifold by $Z_\varphi$. On the boundary this replaces
$S^3_{\rm Hopf}\times D^4$ by $D^4_{\rm Hopf}\times S^3$.
Over $C$, the remaining boundary is exactly the actual free circle
bundle of the clutched sphere: principal circle bundles there are
classified by their Euler class, which is the specified restriction
of $x$. On their common boundary $S^3\times S^2$, a discrepancy
between two such bundle identifications is a circle gauge with a
real logarithm, because $H^1(S^3\times S^2;\Z)=0$.
A collar cutoff corrects it to the prescribed identification.
Alexander duality gives $H^2(C;\Z)=\Z$, generated by the
meridian, so specifying this oriented Hopf restriction fixes the
entire line class on $C$. There is no additional line-bundle
parameter on the exterior. The replacement
$D^4_{\rm Hopf}\times S^3$ is now exactly the
actual fixed tubular neighbourhood with its Hopf normal action.
Hence
\[
 \partial Z_\varphi=\Sigma_\varphi
\]
as an oriented smooth manifold, with the orientation fixed by
agreement on the common free exterior. This is an identification
with the actual boundary, not merely with a sphere having a
conjectured invariant.

Explicitly, the boundary orientation of $D(L_x)$ is the positive
line-circle direction followed by the orientation of its
six-dimensional base. This is the circle-first orientation of
the actual free exterior by the defining choice of $x$.
The surgery product framing then gives the boundary-consistent
orientation on the new fixed tube. If a prior principal-connection
convention uses the dual line, replace $x$ by $-x$ throughout;
both $p$ and $k$ reverse, as does the corresponding integer
conversion. The actual zero evaluation below is unchanged.

The added four-cycle has a particularly concrete description.
Inside the displayed product region, the attachment replaces
\[
 D(\mathcal O(1))\times D^4
 \quad\hbox{by}\quad \mathbb{CP}^2\times D^4.
\]
Indeed a tubular neighbourhood of a projective line in
$\mathbb{CP}^2$ is $D(\mathcal O(1))$, and its complement is a
four-disc, with the same standard Hopf boundary identification.
Let $b$ be this positively oriented $\mathbb{CP}^2\times\{0\}$.
Its normal four-plane bundle is trivial, with the explicit
$D^4$ product framing. In particular $b^2=0$; it is not a
four-sphere with zero Pontryagin evaluation.

Let $a$ be the image of the class $\operatorname{PD}_{V_\varphi}x$
in the zero section of $D(L_x)$. A generic section of $L_x$
represents it by a closed oriented four-submanifold. Its normal
bundle in $Z_0$ is the sum of its normal line in the zero section
and the vertical line, both isomorphic to $L_x$ on that
submanifold. Therefore
\[
 a^2=\langle x^3,[V_\varphi]\rangle=k.
\]
Moreover $a\cdot b=1$: in the surgery product this intersection
is the pairing of $x$ with the meridian $S^2$, which is one by
the Euler-one convention. The four-handle exact sequence and
$H_3(V_\varphi)=0$ give
\[
 0\longrightarrow H_4(V_\varphi;\Z)
 \longrightarrow H_4(Z_\varphi;\Z)
 \longrightarrow\Z\longrightarrow0.
\]
The class $b$ uses the handle core once and maps to its generator.
Thus $(a,b)$ is an integral basis, and its intersection matrix is
the displayed $Q$. Its determinant is $-1$, so it has one
positive and one negative eigenvalue and signature zero.

The handle does not change second cohomology. Let $c$ be the
unique extension of $\pi^*x$ from $Z_0$ to $Z_\varphi$.
On $Z_0$ the tangent bundle restricts stably to
$\pi^*(TV_\varphi\oplus L_x)$, so
\[
 w_2(TZ_0)=\pi^*x\pmod2.
\]
The same degree-two handle isomorphism with $\Z/2$ coefficients
proves $w_2(TZ_\varphi)=c\pmod2$. Thus $c$ is a determinant
class of a spin$^c$ structure. It is trivial on the boundary,
which is a homotopy sphere and has its unique spin structure.
Notice that this coboundary need not be spin: $c$ is primitive.

On $a$ the stable tangent splitting gives
\[
 \langle p_1(TZ_\varphi),a\rangle=p+k,
 \qquad \langle c^2,a\rangle=k.
\]
On $b$ the normal framing gives
$TZ_\varphi|_b=T\mathbb{CP}^2\oplus\R^4$.
The class $c|_b$ is the hyperplane class, since it evaluates as
one on the original meridian projective line. Hence
\[
 \langle p_1(TZ_\varphi),b\rangle=3,
 \qquad \langle c^2,b\rangle=1.
\]
Here $p_1(T\mathbb{CP}^2)=c_1^2-2c_2=3H^2$ follows directly
from its Euler sequence. This proves the stated covectors.

Because the boundary is a homotopy sphere, the relevant relative
and absolute fourth cohomology groups agree. Their characteristic
pairing is the inverse of the integral intersection matrix:
\[
 Q^{-1}=\begin{pmatrix}0&1\\1&-k\end{pmatrix}.
\]
Multiplying the two covectors by this inverse gives exactly
\[
 p_1^2=6p-3k,\qquad p_1c^2=p+k,\qquad c^4=k.
\]

For clarity, the conversion to the boundary invariant uses only
the spin$^c$ index formula, not a knot-classification formula.
The closed eight-dimensional index density is
\[
 [e^{c/2}\widehat A]_8
 =\frac{7p_1^2-4p_2}{5760}
       -\frac{p_1c^2}{192}+\frac{c^4}{384}.
\]
Eliminating $p_2$ by the signature formula and gluing to any spin
filling of the boundary gives \ref{push:src:C.2.1}'s expression
\[
 \mu(\Sigma_\varphi)\equiv
 \frac{p_1^2-4\sigma}{896}
 -\frac{p_1c^2}{192}+\frac{c^4}{384}\pmod\Z.
\]
The primary normalization is
\citet[Section 2.2, equation (2.1)]{CrowleyNordstrom2020},
with its invariant divided by $28$; the spin normalization is
\citet[Section 6, equation (11)]{EellsKuiper1962}.
Substitution now gives, by exact rational arithmetic,
\[
 \frac{6p-3k}{896}-\frac{p+k}{192}+\frac{k}{384}
       =\frac{p-4k}{672}.
\]

Finally \ref{push:src:C.1.1} identifies the mapping torus with
$(S^1\times S^6)\#\Sigma_\varphi$. Boundary connected sum
of $Z_\varphi$ with $S^1\times D^7$, using a trivial determinant
line on the latter, realizes this actual boundary. Either spin
structure on the product can be used. The added summand has no
fourth cohomology and has zero signature; the displayed relative
characteristic numbers do not change.
\end{proof}

\begin{remark}[Scope and actual six-loop evaluation]
For the actual annular six-loop cap the independent geometric
calculations \ref{push:src:C.3.18}--\ref{push:src:C.3.22} give
$p=-24\epsilon_{\rm cap}$ and $k=-6\epsilon_{\rm cap}$.
Here the one common sign includes the possible simultaneous
meridian-sign change just specified.
The present coboundary therefore has
\[
 p_1^2=-126\epsilon_{\rm cap},\qquad
 p_1c^2=-30\epsilon_{\rm cap},\qquad
 c^4=-6\epsilon_{\rm cap},\qquad \sigma=0.
\]
Its spin$^c$ formula is zero directly:
\[
 \epsilon_{\rm cap}
 \left(-\frac{126}{896}+\frac{30}{192}
                         -\frac{6}{384}\right)=0.
\]
This does not use the assertion $J_\psi=0$ or the Wall--Davis
conversion. By itself it evaluates the six-loop sphere. The
independent integer additivity needed for the one-loop conclusion
is supplied separately in \ref{push:src:C.3.26}. The characteristic arithmetic
is checked in \texttt{d\_push\_eight\_coboundary.py}.

This is a smooth topological spin$^c$ coboundary. No complex
structure on it is asserted, and it is not identified with the
specific analytic filling requested in Part B.
\end{remark}
\Uses{the hypotheses, notation and definitions stated here; \cref{push:src:C.3.18}; \cref{push:src:C.3.22}; \cref{push:src:C.3.26}}


